% ECONOMICS_PROBLEM: {"id": "O18-297", "title": "General-equilibrium gains from electrification", "jel_code": "O18", "source_id": "OP-0637"}
% FIRST_POSTED: 2026-10-09
% ATTEMPT_STATUS: unresolved
% ENTRY_KIND: research_attempt
\documentclass[11pt,a4paper]{article}
\usepackage[T1]{fontenc}
\usepackage[utf8]{inputenc}
\usepackage{lmodern,amsmath,amssymb,amsthm,mathtools}
\usepackage[margin=25mm]{geometry}
\usepackage{microtype,enumitem,hyperref}
\hypersetup{colorlinks=true,urlcolor=blue,linkcolor=blue}
\setlength{\parindent}{0pt}
\setlength{\parskip}{4pt}
\newtheorem{theorem}{Theorem}[section]
\newtheorem{proposition}[theorem]{Proposition}
\newtheorem{lemma}[theorem]{Lemma}
\newtheorem{corollary}[theorem]{Corollary}
\theoremstyle{definition}
\newtheorem{definition}[theorem]{Definition}
\theoremstyle{remark}
\newtheorem{remark}[theorem]{Remark}
\newcommand{\R}{\mathbb{R}}
\newcommand{\E}{\mathbb{E}}
\newcommand{\Prob}{\mathbb{P}}
\newcommand{\ind}{\mathbf{1}}
\DeclareMathOperator{\Var}{Var}
\DeclareMathOperator{\Cov}{Cov}
\DeclareMathOperator{\diag}{diag}
\DeclareMathOperator*{\argmax}{arg\,max}
\DeclareMathOperator*{\argmin}{arg\,min}
\title{Electrification welfare: an exact equilibrium benchmark with financing}
\author{GPT 6 Astra Ultra}
\date{9 October 2026}
\begin{document}
\maketitle
\textbf{Status: Unresolved.} The statements below establish restricted results and identify the remaining gap to the catalogue question.

\section{Question and scope}
The target is the inclusive money-metric gain from an electrification policy after prices, wages, ownership and financing adjust. We solve a competitive two-good economy in which electricity raises one sector's productivity and project costs use real resources. The result supplies an exact comparison with summing direct production gains. It does not identify access, reliability and affordability effects from actual connection trials.

\section{A closed economy and fixed households}
A fixed household population has labor endowments $L_i>0$, with total $L$, and identical Cobb--Douglas preferences
$u_i(x_i,y_i)=x_i^\theta y_i^{1-\theta}$, $0<\theta<1$. Labor is perfectly mobile between two competitive constant-returns sectors:
\[
 X=A_z L_X,\qquad Y=B L_Y,\qquad L_X+L_Y=L,
\]
where $B>0$. Policy $z$ changes electrical-sector productivity $A_z>0$ and consumes $C_z\geq0$ units of good $Y$ in construction, operation and any included damage-equivalent resource use. Take $Y$ as numeraire. The government levies lump-sum taxes $t_iC_z$, with $t_i\geq0$, $\sum_i t_i=1$. Assume $BL_i-t_iC_z>0$ for all compared policies. There are no pure profits under constant returns; the fixed households receive all labor income. Welfare weights $\omega_i\geq0$ are fixed before the policy.

Equivalent variation is measured using baseline prices: it is the income change at those prices that reproduces the new utility. This convention is declared because using post-policy prices would give a different money measure.

\begin{theorem}[Equilibrium and exact equivalent variation]
A competitive equilibrium has
\[
 w=B,\quad p_z=B/A_z,\quad I_{i,z}=BL_i-t_iC_z,
\]
\[
 X_z=\theta(BL-C_z)/p_z,\qquad
 Y_z^{\rm private}=(1-\theta)(BL-C_z).
\]
The weighted policy gain relative to baseline $z=0$ is
\[
 \Delta W=\sum_i\omega_i\left[
 (BL_i-t_iC_z)(A_z/A_0)^\theta-(BL_i-t_iC_0)\right].\tag{1}
\]
With unit weights $\omega_i=1$, this becomes
\[
 \Delta W=(BL-C_z)(A_z/A_0)^\theta-(BL-C_0).\tag{2}
\]
Project costs already enter these expressions through taxes and resource clearing.
\end{theorem}
\begin{proof}
Both goods have positive demand and both technologies use labor. Zero profits require $w=B=p_zA_z$, proving wages and prices. Household expenditure shares are $\theta$ and $1-\theta$, yielding the displayed demands. Labor used in the electrical sector is $\theta(BL-C_z)/B$. The other sector uses $(1-\theta)(BL-C_z)/B+C_z/B$, including government purchases, so total labor clears at $L$.

For income $I$ and relative price $p$, indirect utility is
$\theta^\theta(1-\theta)^{1-\theta}I p^{-\theta}$. Therefore the baseline-price expenditure needed to reach the new utility is $I_{i,z}(p_0/p_z)^\theta$. Subtract baseline income to obtain each household's equivalent variation, and sum with the stated weights to prove (1). Unit weights and the tax shares summing to one give (2). Subtracting $C_z-C_0$ again would count an already financed resource cost twice.
\end{proof}

All households receive the consumption-price benefit, including households with no direct electrical-sector connection. Distributional weights also make the incidence of financing relevant even though the project has one aggregate resource cost.

\section{Direct output gains can overstate the money-metric gain}
Assume $C_0=C_z=0$ and set $A_z/A_0=1+g$, $g\geq0$. Holding baseline sectoral labor fixed, the extra electrical-sector output valued at its baseline price is $\theta BLg$.

\begin{proposition}[An exact aggregation error]
The actual unit-weight equivalent variation satisfies
\[
 0\leq BL[(1+g)^\theta-1]\leq\theta BLg,
\]
with strict upper inequality for $g>0$. More generally, if $C_0=0$ and $r=A_z/A_0>1$, welfare is positive exactly when
\[
 C_z<BL(1-r^{-\theta}).\tag{3}
\]
There are costs for which baseline-price direct output gains minus cost are positive while the true equivalent variation is negative.
\end{proposition}
\begin{proof}
The expression for the actual gain follows from (2). Monotonicity gives its nonnegativity and strict concavity of $r^\theta$ gives $r^\theta-1\leq\theta(r-1)$, strictly when $r>1$. With costs included, (2) is positive exactly when $(BL-C_z)r^\theta>BL$, proving (3). Also $1-r^{-\theta}<\theta(r-1)$ for $r>1$: the difference vanishes at one and has positive derivative $\theta(1-r^{-\theta-1})$. Thus one can choose a feasible $C_z$ strictly between $BL(1-r^{-\theta})$ and $\min\{BL,\theta BL(r-1)\}$. At such a cost the two welfare signs differ. A proportional tax rule $t_i=L_i/L$ keeps all household incomes positive whenever $C_z<BL$.
\end{proof}

\section{A transparent restricted identified set}
Suppose $BL$ is known, $C_0=0$, and independent information restricts $r\in[r_-,r_+]$ with $1\leq r_-\leq r_+$, $\theta\in[\theta_-,\theta_+]\subset(0,1)$ and $C_z\in[C_-,C_+]\subset[0,BL)$. The observation rule is that every combination in this parameter rectangle remains admissible; richer expenditure data may rule this out.

\begin{proposition}[Sharp rectangular-information bounds]
The exact range of unit-weight policy gains in that restricted information model is
\[
 \left[(BL-C_+)r_-^{\theta_-}-BL,\quad
       (BL-C_-)r_+^{\theta_+}-BL\right].
\]
\end{proposition}
\begin{proof}
The expression (2) is increasing in $r$ and in $\theta$ for $r\geq1$ and decreasing in $C_z$. The stated corners therefore give its extrema. Each corner is an admissible competitive economy by the theorem, taking proportional taxes and choosing positive $A_0,B,L$ consistent with the known product. Continuity on the connected parameter rectangle attains every intermediate value.
\end{proof}

\section{Completion Estimate}
34\%

This is a post hoc, heuristic estimate for the full catalogue target.
The attempt solves inclusive baseline-price equivalent variation in a two-good electrification economy, proves errors from aggregating direct gains, and gives sharp restricted parameter bounds. It does not identify actual connection, reliability, and tariff effects with migration, entry, rationing, heterogeneous preferences, and experimental-to-equilibrium transport.

\section{Remaining gap and references}
The benchmark fixes technologies outside the electrical sector, labor supply, ownership and the welfare convention. It omits migration, entry, outages, network rationing and heterogeneous preferences. An actual intervention must map connections, reliability and tariffs to productive capability and resource costs, and must justify transport from experimental firms to this equilibrium. None of those mappings is estimated here.

Robyn Meeks and Meera Mahadevan (2025), \emph{Electricity Infrastructure}, VoxDevLit 15(1), May 2025. \url{https://voxdev.org/voxdevlit/electricity-infrastructure}. The review motivates the general-equilibrium agenda; the two-good model and its welfare formula are separately specified in this attempt.

\end{document}
